\section{Scope of the proof of concept}
\label{sec:scope}

\subsection{What the PoC set out to show}

The proof of concept (PoC) had one narrow engineering goal: build a working, checkable chain
from a named domain or capability to a ranked list of questions for a human expert. It used only
public evidence, with zero participants at any step. \evobs{} \emph{Checkable} means three things.
First, every criterion-labeled claim's evidence is an exact, byte-identical span in a document the
system fetched itself. A claim the extractor could not ground this way stays in the ledger as
\code{synthetic\_extrapolation} and never acts as a criterion (\cref{sec:evidence-model}). Second,
every epistemic label is computed from that evidence rather than typed in by hand. Third, every
stage's tests run offline. The chain has four stages:
\begin{enumerate}[noitemsep]
  \item reconstruct a domain's written record with provenance (\cref{sec:architecture});
  \item keep that record as an explicit evidence ledger with six epistemic labels
    (\cref{sec:evidence-model});
  \item apply methodology lenses that ask whether a named construct's content is stated anywhere
    in the ledger; and
  \item rank the resulting candidates into expert questions.
\end{enumerate}
Each stage is implemented code, run on real web sources, with a committed test suite.

The PoC was not meant to show, and does not show, that the resulting questions or their ranking
point at what a real expert actually knows and the record does not. That is a separate claim
(RQ-B, \cref{sec:question}) that needs human evidence to test, and the PoC deliberately used none.
\evint{} What the PoC demonstrates is a mechanism: a chain that turns silence in the written
record into a labeled, ranked hypothesis. Whether that hypothesis is right is untested. The gap
map's own closure step does not yet beat a fair control (\cref{sec:n3,sec:not-demonstrated}), so
even the mechanism's central move is an open question, not a working result.

\subsection{Why zero participants}

The PoC was built entirely inside what \repo{REORIENTATION.md} (\S22) calls the NOW program:
zero participants, zero proprietary or customer data, zero microphones, and zero live expert
interviews. \evobs{} The one human input the program allows at all is a single blind second
coder, used for labeling rather than participation, and that step has not yet been exercised.
The constraint follows directly from the project's own history (\cref{sec:track-a-history}
below): every gate that needed a human to run first stalled on recruiting that human. A design
that needs no participant to reach its first checkable result can be built, tested, and revised
on the calendar of the engineering team alone. Its failures are attributable to the code and
the public evidence, not to a small or biased human sample. \evint{} The cost is symmetric: a
zero-participant PoC cannot, by construction, validate the one claim that matters most --- that
the reconstruction's thin spots correspond to what humans actually know. That validation is
future research (\cref{sec:not-demonstrated,sec:residual,sec:validation}), pre-registered before
any gold data is acquired.

\subsection{How the project reached this design}
\label{sec:track-a-history}

The project began as a physics-education study. Twelve physicists would be interviewed with an
AI-assisted think-aloud and retrospective-probe protocol, and about 370 first-year students would
sit a two-arm randomized trial with delayed transfer. Building the instrument for that study --- a
session app, a coding pipeline, five architecture decisions, and four rounds of lessons drained
into the codebase --- absorbed most of the project's effort. \evobs{} The instrument is complete
and tested: \repo{instrument/} runs 311 tests (\code{uv run --directory instrument pytest -q}).
Every one of its gates, though, from the first pilot session to the final randomized trial,
needed physicists, students, or both to exist before any result could be read.

A self-diagnosis in \repo{REORIENTATION.md} \S3 named this pattern \emph{drift}: a
falsification discipline that demanded validation against observed human performance had,
with no participants recruited yet, nowhere to go except building the instrument ever more
carefully. \evint{} Public-evidence reconstruction, organizational knowledge, and
evidence-grounded synthetic cohorts had never been considered as alternative first steps, so
work of that kind could not register as progress against the existing plan. The owner adopted a
reorientation on 2026-09-28 (\repo{REORIENTATION.md}, \S3, \S22). It freezes the physics study
and its instrument as \emph{Validation Track A}, unchanged and started only on its own registered
preconditions (experts, a course, ethics approval). It also opens a new, parallel program --- NOW,
items N0 through N9 --- whose every item can start and end without a single participant. N0--N3 are
built; \cref{tab:now-verdicts} (\cref{sec:evolution}) tracks what each one did and the verdict it
read. \Cref{tab:now-items} lists the six items that are not yet started.

\begin{table}[tbp]
  \centering
  \small
  \caption{The NOW program's not-yet-started items, N4--N9 (\repo{REORIENTATION.md}, \S22); N0--N3
  are tracked separately in \cref{tab:now-verdicts}. Every item, built or not, is pre-registered to
  end in a continue/change/stop reading of an experiment result, not in a build
  (\cref{sec:evolution}).}
  \label{tab:now-items}
  \begin{tabularx}{\linewidth}{@{}llX@{}}
    \toprule
    Item & Status & One line \\
    \midrule
    N4 & Not started & Behavioral validity: does the reconstructed model fit real learner
      responses? \\
    N5 & Not started & Residual measurement: estimate unseen items against gold as an
      independent occasion. \\
    N6 & Not started & Question selection by expected information gain against held-out gold. \\
    N7 & Not started & Evidence-grounded synthetic cohorts, validated in tiers, never a
      criterion. \\
    N8 & Not started & E-SEAL: sealed Track A predictions, secondary and non-gating. \\
    N9 & Not started & E-COST: a running ledger of cost per validated item. \\
    \bottomrule
  \end{tabularx}
\end{table}

This report covers N0--N3: the four items that exist. \evfut{} N4--N9 are described only as the
research program that would follow (\cref{sec:validation,sec:roadmap}); nothing under those
items has been built or run.
